\section{Ordinary convergence on the large theta side}
\label{sec:sr-large-theta-convergence}

We prove the ordinary-convergence clause of
\(\srcid{ThetaIntegralsAgree.Decay}\).
Let \(E/F\), \(V\), \(W\), the additive character, splitting
characters and rational Schr\"odinger realization be those
of the theta packet.
Write
\[
       [F:\mathbb Q]=f,\qquad
       \dim_E V=a,\qquad \dim_E W=d=a+a'.
\]
The Hermitian signatures of \(V\) are
\((a-p,p)\), with the same \(p\) at every real place.
Let \(r\) be its Witt index over \(F\).
Then \(r\le p\).
The following argument proves convergence whenever
\begin{equation}\label{eq:sr-large-weil-range}
                         d<2(a-r).
\end{equation}
In the source,
\[
     (a-a')/2>2(p+q)+1
\]
implies \(a-a'>2p\ge2r\), hence
\eqref{eq:sr-large-weil-range}.

Fix an adelic Schwartz function \(\Phi\) in the actual
Weil model.  For every \(h\in U(W)(\mathbb A_F)\) we claim
\[
 \int_{U(V)(F)\backslash U(V)(\mathbb A_F)}
                 |\Theta_V(\Phi;g,h)|\,dg<\infty .
\]
The claim uses ordinary integration; no finite-place
regularizer is inserted.

\subsection{A mixed model and its lattice bound}

Choose a rational Witt decomposition
\[
                       V=U\oplus V_0\oplus U^*,
                      \qquad\dim_EU=r .
\]
In the mixed Schr\"odinger model for \(V\otimes_EW\),
the variables belonging to \(U\oplus U^*\) form
\(W(\mathbb A_E)^r\); the remaining variables are
those of a rational Schr\"odinger model for \(V_0\otimes_EW\).
Write their infinite coordinates as
\[
                  (u_1,\dots,u_r,\xi).
\]
Each \(u_i\) has real dimension \(2fd\).
The change from the initially specified rational model
uses rational linear transformations and partial Fourier
transforms.
It preserves the theta functional by rational reindexing
and Poisson summation with selfdual measures.
It preserves the Schwartz class as well.
The normalization of the Weil action fixes its scalar,
so that this is a statement about the specified theta
kernel.

Let \(A\) be the split torus acting by \(e^{t_i}\) on
the \(i\)-th line of \(U\) and by \(e^{-t_i}\) on the
paired line of \(U^*\), with the same real scalar at
each real place of \(F\).
The unitary Schr\"odinger action is
\begin{equation}\label{eq:sr-large-cartan-action}
 \omega(a_{\mathbf t})\phi(u_1,\ldots,u_r,\xi)
   =e^{-fd\sum_i t_i}
       \phi(e^{-t_1}u_1,\ldots,e^{-t_r}u_r,\xi).
\end{equation}
Indeed the real determinant on one \(u_i\)-block is
\(e^{2fd t_i}\), and unitarity supplies its inverse
square root.
The determinant character of the hyperbolic torus is
one.  Thus no splitting-character modulus changes
this exponent.

We record the elementary estimate used to sum these
functions over rational points.
If \(L\) is a fixed full lattice in a real space of
dimension \(D\), and \(M>D\), then
\begin{equation}\label{eq:sr-large-shifted-lattice}
       \sum_{z\in L}(1+\|e^{-t}(z+b)\|)^{-M}
                  \le C_{L,M}e^{Dt}
          \quad(t\ge0,\ b\in\mathbb R^D).
\end{equation}
The number of points in a translated ball of radius
\(R\) is bounded by \(C_L(1+R)^D\), uniformly in its
center: place disjoint balls of a fixed radius around
them and compare volumes.
Divide the sum into dyadic annuli of radii
\(e^t2^j\).  Their contributions are bounded by
\(C_Le^{Dt}2^{j(D-M)}\), a convergent geometric series.
The same estimate with \(t=0\) is uniform in \(b\).

Let \(N\) be the unipotent radical of a rational minimal
parabolic compatible with the chosen Witt flag.
Its mixed-model action has the form
\[
 (\omega(n)\phi)(u,\xi)
     =\psi(Q_n(u,\xi))\,
                   \phi(B_nu,\xi+L_nu),
\]
with unit-modulus character factor and a block unitriangular
\(B_n\).
The block order may be chosen so that its \(i\)-th coordinate
is \(u_i\) plus a linear expression in later coordinates.
This formula follows directly from the three types of
unipotent matrices: changes within \(U\) give \(B_n\),
maps between \(V_0\) and the hyperbolic pairs give
Heisenberg translations and modulations in \(\xi\),
and the central Hermitian block gives the quadratic
character.
Only the linear translation part of each Heisenberg
operator affects absolute values.

If \(n\) ranges over a compact subset of \(N_\infty\),
and the remaining right-hand factors range over compact
sets, all Schwartz seminorms needed below have common
bounds.
For the ordinary Weil representation this follows from
the continuity of linear substitutions and partial
Fourier transforms on the Schwartz space; compact sets
are covered by finitely many such coordinate charts.
At finite places a fixed compactly supported
Schwartz--Bruhat factor restricts rational points to
one fixed fractional lattice and has bounded absolute
value.
Enlarge that lattice to a product lattice in the
displayed rational coordinates.

Apply a product Schwartz bound, then sum first over
\(\xi\).
The \(t=0\) case of
\eqref{eq:sr-large-shifted-lattice} is uniform in its
shift \(L_nu\).
Sum next over \(u_1\), then \(u_2\), and so forth.
The triangular form makes each translation a fixed
shift while that block is being summed.
Equation \eqref{eq:sr-large-shifted-lattice} contributes
\(e^{2fd t_i}\) from the \(i\)-th block.
Together with the prefactor in
\eqref{eq:sr-large-cartan-action}, this proves
\begin{equation}\label{eq:sr-large-theta-bound}
             |\Theta_V(\Phi;n a_{\mathbf t}m k,1)|
                 \le C_\Phi e^{fd\sum_i t_i}
\end{equation}
on each such set, for \(t_i\ge0\).
This summation argument includes the zero rational
vector and all singular orbits.

\subsection{Integration on reduction sets}

We use the usual arithmetic reduction theorem in the
form of finitely many adelic classes and finitely many
Siegel sets for their arithmetic groups.
For a compatible minimal parabolic these sets have
compact unipotent and right-hand factors and split
coordinates
\[
                    t_1\ge t_2\ge\cdots\ge t_r\ge0 ,
\]
after harmless fixed shifts of the origin.
The anisotropic Levi factor has compact arithmetic
quotient.  The norm-one part of its restriction-of-scalars
torus is compact modulo units.
Thus the split coordinates may be taken as the common
\(F\)-scalings used in
\eqref{eq:sr-large-cartan-action}.
No estimate uniform in the field or the reduction
sets is asserted here; ordinary convergence requires
none.

For precision, the positive relative root spaces of
\(U(V)\) are
\[
 \begin{array}{c|c}
        e_i-e_j,\ e_i+e_j\ (i<j)&2\\
        2e_i&1\\
        e_i&2(a-2r).
 \end{array}
\]
The second column gives their dimensions over \(F\).
They can be read off from the off-diagonal
\(E\)-entries, the one-dimensional skew-Hermitian diagonal
entry, and the \(V_0\)-row entries in a Witt matrix.
Summing their weights gives
\[
                  2\rho(\mathbf t)
                      =2\sum_{i=1}^r(a-2i+1)t_i .
\]
Restriction of scalars multiplies this by \(f\).
Consequently Haar measure on the reduction set is
bounded by a constant times
\[
             e^{-2f\sum_i(a-2i+1)t_i}
                           \,dt_1\cdots dt_r
\]
times the finite measures of its compact factors.

Multiply this density by
\eqref{eq:sr-large-theta-bound}.
Put \(t_{r+1}=0\) and \(s_j=t_j-t_{j+1}\ge0\).
The exponent becomes
\[
 \begin{split}
 f\sum_{i=1}^r\{d-2(a-2i+1)\}t_i
        &=f\sum_{j=1}^r
             j(d-2a+2j)s_j .
 \end{split}
\]
Every coefficient is strictly negative by
\eqref{eq:sr-large-weil-range}.
The integral is bounded by the product
\[
                 C_\Phi\prod_{j=1}^r
                    \frac1{fj(2a-d-2j)}<\infty .
\]
When \(r=0\), the arithmetic quotient is compact and
the continuous theta function is integrable directly.
This proves convergence on all reduction sets and hence
on the quotient.

Finally fix any \(h\).
The vector \(\omega(1,h)\Phi\) is Schwartz, so the same
argument applies to it.
The two unitary-group actions commute, and its theta
function at \(g\) is the original kernel at \((g,h)\).
Automorphy under \(U(V)(F)\) follows from rational
reindexing and Poisson summation; the splitting
characters are Hecke characters, hence are trivial on
rational determinants.
Integrability is therefore independent of the
fundamental domain.
Any Haar measure differs by a positive scalar, which
does not affect finiteness.
We have proved exactly
\[
          \forall\Phi\in\mathcal S,\qquad
               \srcid{Model.Conv}(M,\omega,\Phi).
\]

This proves the large-side ordinary-convergence
assertion of the packet. The next subsections identify
the prescribed split bar test with the global rank cutoff
and establish the small-side clauses separately.



\subsection{Matching the specified bar test and the global coefficient cutoff}
\label{subsec:sr-exact-bar-cutoff}

We verify that the local test used above is precisely the
source's \(\mathrm{SoLoc}\), including its existential
choice of split intertwiner. Let \(v_1\) be the designated
split finite place, \(k=F_{v_1}\), and \(t:E\to k\) its
chosen embedding. Write
\[
 P=J_V(t),\qquad Q=J_W(t),\qquad
             \operatorname{coB}(z)=-Pz_{t\sigma}Q^{\mathsf t}.
\]
Both \(P,Q\) are invertible. Thus
\(\operatorname{rank}\operatorname{coB}(z)
 =\operatorname{rank}z_{t\sigma}\).
The local condition supplies a Schwartz bijection
\(\iota\) satisfying the split bar-model intertwining
laws, with both auxiliary scalar normalizations equal
to one, and such that \(\iota\phi\) vanishes off the
full-row-rank locus. No special choice of this
\(\iota\) will be made.

Put \(T=\operatorname{supp}(\iota\phi)\).
It is compact and full rank. Since \(2\ne0\) in \(k\),
\(2T\) is full rank as well. Choose a compact open
additive subgroup \(\Lambda\) of \(M_{n,d}(k)\) with
\(2T+\Lambda\) full rank, and let \(K\) be its annihilator
under \((Y,U)\mapsto\psi_v(\operatorname{tr}(YU^{\mathsf t}))\).
One may shrink \(\Lambda\) so that \(T\) meets only
finitely many of its cosets; compactness in the
discrete quotient gives this finiteness for every
compact open \(\Lambda\).
Nondegeneracy of the trace pairing and local Fourier
duality give \(K^\perp=\Lambda\).
The explicit scaled-integral-lattice construction in
Section~\ref{sec:theta-decay-proof} gives the same
annihilator if the additive character has nontrivial
conductor.

Partition \(\iota\phi\) by these cosets and pull the
pieces back by \(\iota^{-1}\). They are Schwartz and
sum to \(\phi\), since \(\iota\) is a linear bijection
on Schwartz functions. Fix a piece with image supported
in \(Y_0+\Lambda\). For \(U\in K\), the defining
bar-model multiplication law says that its local
Schr\"odinger translate by the split vector
\(w_U=\operatorname{spM}(-U,0)\) is
\[
       S(M.\operatorname{co}(w_U))\phi
          =\psi_v(\operatorname{tr}(Y_0U^{\mathsf t}))\phi.
\]
This is the \(\mathrm{IsLinBar.mult}\) law itself;
it does not require \(\iota\) to be an isometry.

The symplectic pairing in the specified rational model is
\[
 \langle z,w_U\rangle
                 =\operatorname{tr}(\operatorname{coB}(z)U^{\mathsf t}).
\]
For the reflected coefficient \(c_\Phi\), the elementary
Schr\"odinger commutator calculation therefore gives
\[
 c_\Phi(M.\operatorname{co}(z))
   =\psi_v\!\left(
        \operatorname{tr}((2Y_0-\operatorname{coB}(z))U^{\mathsf t})
             \right)c_\Phi(M.\operatorname{co}(z)).
\]
If \(z_{t\sigma}\) has rank less than \(n\), then
\(\operatorname{coB}(z)\notin2Y_0+\Lambda\).
An element \(U\in K\) makes the displayed scalar
different from one. The coefficient is consequently
zero. This proves the cutoff for each piece and every
choice of the other tensor factors, including the
archimedean Schwartz factor.

Here is why the cutoff transfers to the global rational
sum. The model's rational symplectic-coordinate map is
an injective \(F\)-linear map from \(M_{n,d}(E)\) into
\(F^{nd}\oplus F^{nd}\). Both spaces have dimension
\(2nd\) over \(F\), so it is a bijection.
The coefficient sum can therefore be indexed by
\(x\in M_{n,d}(E)\). For every \(g,h\), a matrix
\(x\) of deficient \(E\)-rank stays deficient after
extension by \(t\sigma\) and the invertible left and
right actions. The preceding local coefficient is
zero at \(g^{-1}x(h^{-1})^{\mathsf t}\).
The reflected theta sum equals the original theta sum:
reflection is the central \((-1,1)\) Weil operator
up to a unit scalar, and rational reindexing sends
\(x\) to \(-x\). The coefficient identity thus bounds
an absolute square. These are exactly the parity and
square identities proved in
Subsection~\ref{subsec:sr-theta-coefficient-proof}.
Taking square roots, adding the finitely many piece
bounds, and enlarging their compact supports proves
\eqref{eq:theta-lattice-majorant} with rational full
row rank for every element of the specified
\(\mathcal S^\circ\).
This matches the source's \(\mathrm{soCut}\) and
\(\mathrm{sqSo}\) constructions without imposing
a more restrictive local test.

\subsection{The rational-translate height bound}
\label{subsec:sr-rational-torus-height}
We also spell out the height assertion used when the
second variable is allowed to vary. Its rational
translate is essential. It cannot be inferred from
submultiplicativity after simply discarding the
rational translating matrix.

In a fixed rational Witt basis put
\(a_{\mathbf v}\) diagonal with weights
\(v_1,\ldots,v_y,-v_1,\ldots,-v_y,0,\ldots,0\)
at all infinite places and identity at finite places.
Write \(V=\sum_jv_j\), with
\(v_1\ge\cdots\ge v_y\ge0\).
For the multiplicative adelic matrix height
\[
 H_0(b)=\prod_w
        \max(1,\|b_w\|_w,\|b_w^{-1}\|_w)^{n_w},
\]
use the sup matrix norm at finite places and the
Euclidean operator norm at infinity; \(n_w\) is the
local degree in the product formula. At all but
finitely many places the local factor is one.
For every \(\delta\in\mathrm{GL}_d(E)\), if \(y>0\),
\begin{equation}\label{eq:sr-height-rational-diagonal}
                   H_0(\delta a_{\mathbf v})
                                      \ge e^{D v_1},
                      \qquad D=[E:\mathbb Q].
\end{equation}
Indeed choose a nonzero entry \(\alpha\) in the row
of \(\delta^{-1}\) corresponding to weight \(-v_1\).
In \((\delta a_{\mathbf v})^{-1}
 =a_{\mathbf v}^{-1}\delta^{-1}\), that entry is
\(e^{v_1}\alpha\) at every infinite place and \(\alpha\)
at every finite place. Each matrix-norm factor is at
least the absolute value of this entry. Multiplication
over all places and the product formula for
\(\alpha\ne0\) give \eqref{eq:sr-height-rational-diagonal}.
A fixed rational basis change contributes only a
fixed multiplicative constant, by submultiplicativity
applied on both sides.

The matrix height \(H\) in the Weil model and \(H_0\)
are polynomially comparable, with constants depending
on the fixed arithmetic model. To verify this directly,
expand the \(E\)-matrix entries in a fixed \(F\)-basis
and compare their component norms with the finitely
many embedding norms at each place.
Finite-dimensional norm comparison gives local
constants in each direction; outside the finitely
many places where this basis, its inverse or the
different is not integral, those constants are one.
Taking their product gives a finite global constant.
Replacing component maxima by their local-degree
products changes the comparison by at most a fixed
integer power. Apply the same calculation to inverse
matrices. Thus \(H_0\le C H^b\) for fixed \(C,b\),
and conversely after changing them.
Since \(V\le yv_1\), \eqref{eq:sr-height-rational-diagonal}
gives fixed \(C_1,b_1\) such that
\[
      e^V\le C_1 H(\delta a_{\mathbf v})^{b_1}
                                  \quad\hbox{for every rational }\delta.
\]
For \(y=0\) this holds with \(e^V=1\).
The upper bound \(H(a_{\mathbf v})\le C_2e^{b_2V}\)
follows from its finitely many eigenweights and the
same fixed-basis comparisons.

Now write a reduced second variable as
\(\gamma h=a_{\mathbf v}\kappa\), with \(\kappa\)
in the fixed compact reduction set.
Use the preceding lower bound with
\(\delta=\gamma^{-1}\).
Since \(\gamma^{-1}a_{\mathbf v}=h\kappa^{-1}\),
compact boundedness, inversion invariance and
submultiplicativity give
\begin{equation}\label{eq:sr-height-second-variable}
             e^V\le C_3H(h)^{b_1}.
\end{equation}
This supplies the exact height step in the
two-chamber proof of
\eqref{eq:theta-full-decay}, including all its
rational translates.



\subsection{Full-rank tests on the small theta side}
\label{subsec:sr-small-theta-ordinary-convergence}

The full-rank test condition also gives ordinary absolute
convergence when the first space is the smaller member.
For this assertion the second adelic variable is fixed;
a height estimate uniform in that variable is unnecessary.

Let \(V\) have \(E\)-dimension \(n\), let \(W\) have
\(E\)-dimension \(d\ge n\), and use the actual rational
Weil model and the full-rank subspace \(\mathcal S^\circ\)
at its distinguished split finite place. For every
\(\Phi\in\mathcal S^\circ\) and every fixed
\(h\in U(W)(\mathbb A_F)\),
\begin{equation}\label{eq:sr-small-theta-absolute}
 \int_{U(V)(F)\backslash U(V)(\mathbb A_F)}
             |\Theta(\Phi;g,h)|\,dg<\infty .
\end{equation}
No dimension inequality beyond \(n\le d\) is needed.

Here are the details of the cusp estimate. The coefficient
identity and compact-coset separation in
Subsection~\ref{subsec:sr-theta-coefficient-proof} give, for
each exponent \(L\), the full-row-rank majorant
\[
 |\Theta_{f\otimes\phi}(g,1)|^2
 \le p_{\phi,L}(f)^2
       \sum_{\substack{x\in M_{n,d}(E)\\ \operatorname{rank}x=n}}
       {\bf1}_{C_f}((g^{-1}x)_f)
                 (1+|g^{-1}x|_\infty)^{-L}.
\]
Its continuous seminorm and finite compact support are
chosen before \(f\) and \(g\). In particular the rank
condition concerns the coefficient matrix, after taking
the absolute square; it is not an assertion that the
original theta summation variable has full row rank.

Put \(D=[E:\mathbb Q]\), and let \(r\) be the rational
Witt index of \(V\). Reduction theory gives finitely
many rational Witt frames and representatives
\[
       \gamma g=a_{\mathbf u}\kappa,\qquad
       u_1\ge\cdots\ge u_r\ge0,\qquad \kappa\in\Omega ,
\]
with \(\Omega\) compact. The archimedean eigenweights
in such a frame are \(u_1,\ldots,u_r\),
\(-u_1,\ldots,-u_r\), and \(n-2r\) zeros;
the torus is the identity at every finite place.
Rational invariance reindexes the matrix sum.
The inverse matrices of \(\kappa\) have uniformly
bounded archimedean operator norms, and their finite
components enlarge \(C_f\) to one fixed compact set.
A fixed rational change to the Witt basis therefore
leaves only matrices
\(\xi\in M_{n,d}(m^{-1}\mathcal O_E)\)
for one integer \(m\ge1\). Up to fixed positive
constants their infinite norm is
\[
 s_{\mathbf u}(\xi)=
       \sum_{v,i,j}e^{-\lambda_i(\mathbf u)}
                                    |\tau_v(\xi_{ij})|.
\]
All constants in this paragraph can be taken
simultaneously over the finite frame list.

For \(L_0=(D+1)nd\), the one-coordinate lattice packing
argument gives
\begin{equation}\label{eq:sr-small-all-lattice}
 \sum_{\xi\in M_{n,d}(m^{-1}\mathcal O_E)}
                   (1+s_{\mathbf u}(\xi))^{-L_0}
                    \le C\exp\!\left(Dd\sum_i u_i\right).
\end{equation}
To see the exponent directly, use
\[
 (1+\sum_{i,j}A_{ij})^{-L_0}
            \le\prod_{i,j}(1+A_{ij})^{-(D+1)}
\]
with \(A_{ij}=e^{-\lambda_i}\sum_v|\tau_v(\xi_{ij})|\),
and enlarge to independent integral-ideal lattice entries.
Packing in real dimension \(D\) bounds each factor sum
by \(C_m e^{D\max(0,\lambda_i)}\).
There are \(d\) entries per row and the positive
eigenweights sum to \(\sum_i u_i\), giving
\eqref{eq:sr-small-all-lattice}.

If \(r>0\), every full-row-rank \(\xi\) has a nonzero
entry \(\alpha\) in its row with weight \(-u_1\).
The algebraic integer \(m\alpha\ne0\) has absolute norm
at least one. Consequently
\(\max_v|\tau_v(\alpha)|\ge m^{-1}\), and
\[
              1+s_{\mathbf u}(\xi)\ge m^{-1}e^{u_1}.
\]
Write \(U=\sum_i u_i\le r u_1\).
For every integer \(q\ge0\), take
\(M=Ddr+2qr\). Multiplying the preceding estimate to
the power \(M\), and then using
\eqref{eq:sr-small-all-lattice}, gives
\[
 \begin{split}
 \sum_{\operatorname{rank}\xi=n}
       (1+s_{\mathbf u}(\xi))^{-(L_0+M)}
 &\le m^M C e^{DdU-Mu_1}\\
 &\le m^M C e^{-2qU}.
 \end{split}
\]
The square majorant therefore proves
\[
           |\Theta_{f\otimes\phi}(g,1)|
                        \le C_{\phi,q}p_{\phi,q}(f)e^{-qU}
\]
on each reduction set. The finite constants can be
absorbed into a continuous seminorm. When \(r=0\),
the same conclusion with \(U=0\) follows from the
unrestricted lattice bound on the compact reduction set.

In particular \(q=0\) bounds the theta function uniformly
in \(g\). The rational unitary quotient has finite
Haar volume, so this bound proves its absolute
integrability. More generally, the upper torus height
bound gives \(H_Q(g)\le A e^{bU}\). Taking \(q\ge bR\)
gives rapid decay against \(H_Q(g)^R\), for every
integer \(R\ge0\), with the second variable still fixed.

Finally the two Weil actions commute. The action of
\((1,h)\) preserves \(\mathcal S^\circ\), by the
split-matrix rank stability proved in
Section~\ref{sec:sr-good-split-coinvariants}.
Replace \(\Phi\) by \(\omega(1,h)\Phi\) in the preceding
argument. Its archimedean action is continuous on
Schwartz space, and its finite action produces a finite
sum of the allowed pure tensors. This proves
\eqref{eq:sr-small-theta-absolute} for every fixed \(h\)
and every \(\Phi\in\mathcal S^\circ\).
The same bound is uniform for \(h\) in a compact set:
the Schwartz seminorms have bounded operator constants
there and the finite factors have common compact support
and a common local constancy level.

This establishes ordinary convergence on the small side
with its actual full-rank test condition.
The bounds here allow constants depending on the fixed
arithmetic model and on \(h\); they do not assert
a uniform metric or covolume estimate over the varying
family, nor the global equality of the two integrals.

\subsection{The complete ordinary-convergence and decay clause}
\label{subsec:sr-theta-decay-closed}
For every actual triple in the packet,
the large-side ordinary convergence proved above,
the just-matched full-rank coefficient cutoff, and
the two-chamber argument of
Section~\ref{sec:theta-decay-proof} give the following.
For every small-side \(\Phi'\in\mathcal S^\circ\)
and every integer \(R\ge0\), there are \(C\ge0,k\ge0\)
such that
\[
       |\Theta'(\Phi';g,h)|H_Q(g)^R\le C H(h)^k
                                      \quad\hbox{for all }g,h .
\]
The exponent in that chamber proof is chosen before
the finite test factor; all seminorms remain continuous.
Finite sums of pure tensors are handled by the
triangle inequality and by replacing their exponents
with the maximum, since \(H(h)\ge1\).
At \(R=0\), integration over the finite-volume first
quotient gives
\[
 \int |\Theta'(\Phi';g,h)|\,dg\le C' H(h)^k,\qquad
                  |I'(\Phi';h)|\le C'H(h)^k .
\]
The integrand is continuous, hence measurable.
These inequalities prove ordinary absolute convergence
and exactly the moderate-growth assertion for \(I'\).
They use its specified Haar measure; a fixed positive
normalizing scalar only changes \(C'\).

Together with the unrestricted large-side convergence,
these are all clauses of the source's \(\mathrm{Decay}\)
predicate. They hold for the actual bar-test subspace,
rather than a substitute subspace. They do not
establish \(\mathrm{OneConst}\), the two rank recursions,
or the covolume/metric bound; those assertions have
different quantifiers and content.
